% 
%  chapter1.tex
%  ISELthesis
%  
%  Created by Matilde Pós-de-Mina Pato on 2012/10/09.
%  Revised in 2026 (version 4.3.2)
%
\chapter{Introduction}
\label{cha:introduction}

This document is the \LaTeX{} template for bachelor's project reports and master's dissertations at the Lisbon School of Engineering (ISEL), Polytechnic University of Lisbon. It follows the ISEL style guide \emph{Normas de escrita e apresentação de trabalhos finais de curso}, and it is also its own user manual: the PDF you are reading was produced by the template, and its source files are examples you can learn from and replace with your own text.

\section{Before You Start}
\label{sec:before_you_start}

\noindent%
\fcolorbox{iselorange}{white}{%
\begin{minipage}{\dimexpr\linewidth-2\fboxsep-2\fboxrule}
\begin{enumerate}
	\item \textbf{Compile with LuaLaTeX}, not pdfLaTeX, so that the Verdana font required by ISEL is used (Chapter~\ref{cha:a_short_latex_tutorial_with_examples}).
	\item \textbf{Edit only} the \texttt{Config} and \texttt{Chapters} folders and your bibliography file (Chapter~\ref{cha:users_manual}).
	\item \textbf{Do not change} \texttt{template.tex}, \texttt{iselthesis.cls} or the \texttt{ISELthesis-files} folder.
\end{enumerate}
\end{minipage}%
}

\section{Organisation of this Document}
\label{sec:organisation}

\begin{itemize}
	\item \textbf{Chapter~\ref{cha:users_manual}} describes the folders and files of the template, the document options and how to configure your document.
	\item \textbf{Chapter~\ref{cha:a_short_latex_tutorial_with_examples}} explains how to compile the document and shows examples of common \LaTeX{} constructs: figures, tables, glossaries, algorithms, code listings and citations.
\end{itemize}

\section{License and Status}
\label{sec:license}

This template is distributed under the \href{https://opensource.org/licenses/MIT}{MIT License}.

\begin{center}
	\textbf{\large This package and template are not official for ISEL/UPL.}
\end{center}

Bugs and suggestions are welcome at \url{https://github.com/matpato/iselthesis}.